% What is PM
Precision medicine (PM) aims to improve healthcare decision-making by tailoring medical treatments to individual characteristics \citep{kosorok2015adaptive, 
jiang_estimation_2017,
Moodie2020, 
kosorok_introduction_2021}.
% ITRs
One way to formalize this goal is through individualized treatment regimes (ITRs), maps from patient characteristics to treatment recommendations \citep{kosorok_precision_2019}. Ideally, we would like to identify optimal ITRs, the treatment regimes that, if followed, would yield the best outcome on average in the targeted population \citep{kosorok_precision_2019, kosorok2015adaptive}. One way to evaluate this is through the value function, the expected clinical outcome based on a specified treatment regime.

% Best outcome 
In healthcare settings, identifying the best outcome can be challenging, especially when managing complex health conditions. For example, chronic limb threatening ischemia (CLTI), the most severe form of peripheral artery disease (PAD), is a disease characterized by abnormal narrowing and occlusion of the arteries in the lower extremities. CLTI is a marker for severe, poorly controlled cardiovascular disease and diabetes. Even with medical and surgical treatment, patients with CLTI are at high risk for devastating outcomes, including limb loss, cardiac events (e.g., heart attack and stroke), and death \citep{criqui_lower_2021}. From a statistical perspective, CLTI and CLTI outcomes demonstrate an example of competing risks (CR). Competing risks arise when individuals are at risk for multiple types of events, and the occurrence of one event precludes or alters the probability of other events occurring. For example, if a patient with CLTI dies, they are no longer at risk for limb loss. Standard survival analyses ignore CR by including patients experiencing the competing risk as part of the risk set, despite that they can no longer experience the main event of interest. Thus, methods estimating the cumulative incidence function (CIF), or subdistribution, is often of interest when analyzing CR data \citep{kalbfleisch1980statistical, fine_proportional_1999, austin_introduction_2016}. By accounting for competing risks, estimation methods avoid inflation or bias in event probabilities or survival estimates.

Patients with CLTI are highly heterogeneous, and evidence suggests that precision medicine strategies may improve treatment recommendations and, ultimately, patient outcomes \citep{mcginigle_precision_2021, chung_analysis_2022}. Current ITR estimation methods for handling the competing risk setting, like that posed by CLTI, focus on either optimizing the survival curve or minimizing risks of individual events, i.e., minimizing the CIF. However, none of the existing methods \textit{simultaneously} optimize the survival function and the CIF. 

% Our Contributions
We address this gap by proposing a nonparametric two-phase optimal ITR estimator that leverages multi-utility optimization. To do this, we construct a novel multi-utility value (objective) function that maximizes overall survival (OS) from all causes while minimizing the CIF of a priority cause (PC), the main event of interest over any other cause of failure. The first phase of our estimation strategy constructs individual OS curves by adapting the generalized random survival forests from \cite{cho_multi-stage_2023} to the single stage treatment setting. We develop random cumulative incidence forests (RCIFs) to construct individualized CIF curves for the PC in the second phase. To select an optimal ITR, we use a flexible, user-specified criterion (i.e., a tolerance level) to balance treatments so that they appropriately weigh survival and risks. Unlike most existing methods restricted to binary treatments, our method considers a flexible finite number of treatment arms. Because our forests are nonparametric, we can capture complex decision rules. \ifnum\blind=1\else The software package implementing our proposed method is made freely available at https://github.com/cwzhou/itrSurv.\fi

% Paper structure
% Paper structure
The rest of this paper is organized as follows: Section \ref{sec:background} reviews existing literature for precision medicine and the competing risk setting. Section \ref{sec:meth} details our analytic framework and introduces our proposed method. Section \ref{sec:theory} presents the theoretical properties of our ITR estimator and forests. Section \ref{sec:sims} describes the finite sample performance of our proposed estimator through simulation studies. We demonstrate our proposed method's application to the CLTI setting in Section \ref{sec:realdata} and follow with a discussion in Section \ref{sec:disc}.